\documentclass[12pt, a4paper]{article}

% ---------------------------------------------------------------
% NOTE ABOUT INPUTS
% ---------------------------------------------------------------
% The requested template path docs/reference.tex was not found in the codebase.
% The available reference template used for formatting analysis is:
%   docs/reference_report.tex
%
% The requested image folder assets/images was not found in the codebase.
% Project images are available under:
%   assets/
% and are referenced from this report accordingly.

% ---------------------------------------------------------------
% Packages
% ---------------------------------------------------------------
\usepackage[utf8]{inputenc}
\usepackage[T1]{fontenc}
\usepackage{amsmath, amssymb}
\usepackage{graphicx}
\usepackage{booktabs}
\usepackage{array}
\renewcommand{\arraystretch}{1.3}
\usepackage{svg}
\usepackage{enumitem}
\usepackage{geometry}
\usepackage{caption}
\usepackage{subcaption}
\usepackage{cite}
\usepackage{url}
\usepackage{parskip}
\usepackage{longtable}
\usepackage{multirow}
\usepackage{float}

\geometry{
  top=2.5cm,
  bottom=2.5cm,
  left=3.5cm,
  right=3.5cm
}

\hypersetup{
  colorlinks=true,
  linkcolor=blue,
  citecolor=blue,
  urlcolor=blue
}

% ---------------------------------------------------------------
% Document
% ---------------------------------------------------------------
\begin{document}

% ---------------------------------------------------------------
% Title & Author
% ---------------------------------------------------------------
\begin{center}
  {\Large \textbf{Multi-Agent LLM Research Automation Platform: Architecture, Implementation, and System Evaluation}}\\[1em]
  Om Choksi\\[0.25em]
  Devang Dhandhukiya\\[0.5em]
  \small
  Department of Artificial Intelligence and Machine Learning,\\
  Charotar University of Science and Technology, Anand (Gujarat), India\\[0.3em]
  \texttt{23aiml010@charusat.edu.in}, \texttt{23aiml014@charusat.edu.in}
\end{center}

% ---------------------------------------------------------------
% Abstract
% ---------------------------------------------------------------
\begin{abstract}
This report presents a Multi-Agent LLM Research Automation Platform that combines Python,
Node.js, Next.js, PostgreSQL, and Redis across containerized services. The system orchestrates
research through clearly defined stages—discovery, retrieval, synthesis, critique, and
reporting—enabling reliable automation of complex investigative work. By integrating
production patterns such as workspace isolation, health monitoring, and event-driven
traceability, the platform demonstrates how thoughtful architecture supports both performance
and maintainability. Every technical claim in this report maps to verified implementation
artifacts in the codebase.

\textbf{Keywords:} Multi-Agent Systems, LangGraph, LLM Orchestration, FastAPI, Docker.
\end{abstract}

% ---------------------------------------------------------------
\section{Introduction}
% ---------------------------------------------------------------

Research workflows tend to be painfully manual. You jump between tools to find relevant sources, organize references, check whether claims hold up, pull out key insights, and finally write up your findings. Each step usually involves its own verification process, and nothing connects together smoothly. The Multi-Agent LLM Research Automation Platform changes that by coordinating specialized agents across well-defined stages—so you get both speed and a clear trail of what happened. Each stage works on its own but passes results to the next, turning what would otherwise be a brittle collection of scripts into a coherent, trackable workflow.

Under the hood, the system runs five core services inside Docker. PostgreSQL handles long-term data storage. Redis manages temporary state and caching. The AI engine orchestrates all the agent logic and model calls. The backend API takes care of authentication, session management, and routing. And the frontend gives users a clean interface to work with. There's also a dedicated worker that processes jobs asynchronously, so the interface never freezes up even during deep research runs. 

Figure~\ref{fig:system_architecture} (on the next page) shows how these services fit together.

\newpage
\begin{figure}[H]
  \centering
  \includegraphics[width=0.92\textwidth]{assets/system_architecture.png}
  \caption{Multi-service deployment topology with PostgreSQL, Redis, FastAPI engine, Express backend, and Next.js frontend.}
  \label{fig:system_architecture}
\end{figure}

The codebase reflects real‑world engineering practices. Health checks make sure services start in the right order and stay available. Pipeline execution emits events that are stored in the database and streamed to connected clients, so you can track progress live. Search queries fall back gracefully when the first attempt fails. Model providers rotate API keys and support multiple execution modes—from local inference to cloud APIs. Database migrations handle schema changes safely. All these details show that the project cares as much about operational reliability as it does about algorithmic quality.

The platform's agent design follows the same philosophy. Instead of throwing everything into one general‑purpose agent, the system uses specialized agents for discovery, validation, synthesis, and reporting. That way each agent becomes an expert at its own job while staying simple and easy to debug. Figure~\ref{fig:agent_architecture} illustrates how these agents fit into the overall workflow.
\begin{figure}[H]
  \centering
  \includegraphics[width=0.92\textwidth]{assets/agent_architecture_diagram.png}
  \caption{Research pipeline stages: discovery, retrieval, synthesis, validation, and reporting orchestrated through LangGraph.}
  \label{fig:agent_architecture}
\end{figure}

\subsection{Motivation and Problem Context}

Research inherently involves multiple phases: finding relevant literature, analyzing credibility,
extracting key insights, and synthesizing findings into coherent arguments. Existing tools
optimize individual stages in isolation, leaving researchers to manually coordinate work
across systems. The platform eliminates this friction by creating a unified workflow pipeline.
Each stage processes inputs and hands results to the next stage. This approach makes workflows
trackable (every decision is logged), reviewable (intermediate outputs remain visible), and
recoverable (failed stages can be investigated without restarting the entire job).

The implementation achieves this through deliberate architectural choices. Search logic lives
in provider modules. Quality scoring lives in dedicated validation agents. Writing lives in
reporting agents. Infrastructure lives in the backend. This separation makes the system
easier to test, easier to understand, and easier to change.

\subsection{Project Scope}

The repository comprises four primary components:

\begin{itemize}
  \item A Python FastAPI service that orchestrates agent execution and abstracts across
        multiple LLM providers.
  \item A Node.js/Express backend that manages user authentication, session lifecycle, and
        event routing.
  \item A Next.js frontend that presents research workspaces and visualizes execution progress
        in real time.
  \item A PostgreSQL schema with supporting migrations, comprehensive API documentation,
        and architectural notes in the \texttt{docs/} directory.
\end{itemize}

This report focuses exclusively on implementation details verifiable through the actual code,
schema, configurations, and visual assets.

\subsection{Key Technical Contributions}

This implementation demonstrates five important patterns relevant to production AI systems:

\begin{enumerate}[label=\textbf{\alph*}.]
  \item \textbf{Modular Agent-Based Research Orchestration:} Rather than training a single
        generalist model, this system coordinates specialized agents designed for specific
        tasks: discovery, synthesis, validation, and reporting. This separation simplifies
        debugging and enables each agent to develop domain expertise.
  \item \textbf{Provider-Agnostic LLM Runtime:} A factory pattern routes requests across
        Groq, OpenRouter, Google Gemini, local Ollama, and Hugging Face models. Fallback
        logic ensures the system remains functional even when preferred providers become
        unavailable. This flexibility proves critical for production deployments.
  \item \textbf{Workspace-Aware Data Isolation:} Every session and data source links to a
        workspace ID, enabling secure multi-tenant support. Multiple users can conduct
        independent research without interference.
  \item \textbf{Event-Driven Execution Traceability:} Execution events persist in dedicated
        tables and expose through REST APIs, enabling real-time progress visualization on
        the frontend and detailed postmortem analysis for failures.
  \item \textbf{Containerized Multi-Service Deployment:} Docker Compose orchestrates six
        services with health checks, ordered startup, and shared storage. This standardizes
        deployment across development and production environments.
\end{enumerate}

% ---------------------------------------------------------------
\section{Literature Review}
% ---------------------------------------------------------------

This section contextualizes the implementation against distributed AI systems practices.

\subsection{From Single-Agent Systems to Coordinated Pipelines}

Simple language models excel at isolated tasks. Research, however, requires iterative evidence gathering: discovering a topic, retrieving literature, assessing credibility, and synthesizing insights. Single models struggle with such complexity.

This platform separates concerns: topic discovery from retrieval, retrieval from synthesis. Each stage feeds the next, creating clear pipelines and enabling failure recovery.

How this platform implements the pattern is worth noting. The codebase includes agent
inventory files and test endpoints that invite explicit verification at each stage. This
design choice supports both learning and production use: students can study how individual
agents work, while production operators can debug failures by examining stage-level outputs
rather than trying to reason about an end-to-end black box.

\subsection{Orchestration Frameworks and Execution Checkpointing}

Coordinating multiple agents requires balancing flexible control flow with reproducible execution. This platform uses LangGraph with checkpointing, enabling pausable/resumable execution. FastAPI wraps orchestration cleanly: external clients see stable endpoints while internal complexity remains hidden.

\subsection{Heterogeneous Source Retrieval and Quality Assessment}

Research relies on diverse sources with varying schemas and quality levels. This platform queries multiple sources in parallel (web, Wikipedia, arXiv) and normalizes results into standard format: title, URL, description, source ID, favicon, optional dates/thumbnails. Query fallback reduces failed retrievals by trying simplified variants when exact matches fail.

\subsection{Operational Reliability in AI Systems}

Production systems face reliability demands: availability during provider failures, visibility for debugging, zero-downtime evolution, cost management, scaling. This codebase includes rate limiting, health checks, retry counters, workspace isolation, event logging, and migration scripts—foundational infrastructure patterns revealing production maturity.

In summary, the implementation is architecturally sound. It applies modular agent design,
clear service boundaries, event-driven observability, and environment-based configuration.
These principles appear throughout established distributed systems literature and prove
their value here.

% ---------------------------------------------------------------
\section{Materials and Method}
% ---------------------------------------------------------------

\subsection{Codebase Architecture}

The system spans five core components: AI orchestration, backend services, frontend, documentation, and assets.
Table~\ref{tab:repo_structure} summarizes the repository structure.

\begin{table}[H]
\centering
\caption{Repository component summary showing directory roles, primary responsibilities, and technical composition.}
\label{tab:repo_structure}
\begin{tabular}{p{2.4cm}|p{3.6cm}|p{5.0cm}}
\toprule
\textbf{Directory} & \textbf{Primary Role} & \textbf{Notable Contents} \\
\midrule
\texttt{ai\_engine/} & AI orchestration runtime & Agents, graph pipeline, LLM providers, search service, FastAPI app \\
\texttt{backend/} & API gateway and job control & Express routes, worker service, auth middleware, DB integration \\
\texttt{frontend/} & Web client interface & Next.js app, components, charts, markdown rendering \\
\texttt{mobile-app/} & Mobile app scaffold & Flutter structure (future development) \\
\texttt{docs/} & Technical documentation & Architecture notes, API docs, schema, plans, testing guides \\
\texttt{assets/} & Visual artifacts & Architecture diagrams, UI screenshots, workflow visuals \\
\texttt{output/} & Generated outputs & JSON research results and generated report artifacts \\
\bottomrule
\end{tabular}
\end{table}

\begin{figure}[H]
  \centering
  \includegraphics[width=0.92\textwidth]{assets/research_workspace.png}
  \caption{Workspace dashboard with project navigation and session history management.}
  \label{fig:research_workspace}
\end{figure}

\subsection{Deployment Services and Network Topology}

Docker Compose coordinates six services: PostgreSQL, Redis, the Python AI engine, a Node.js backend, a worker process, and a Next.js frontend.
Table~\ref{tab:service_topology} details the topology.

\begin{table}[H]
\centering
\caption{Containerized service topology including network interfaces, port assignments, and technology stack for each component.}
\label{tab:service_topology}
\begin{tabular}{p{1.8cm}|p{2.1cm}|p{3.5cm}|p{6.2cm}}
\toprule
\textbf{Service} & \textbf{Port} & \textbf{Stack} & \textbf{Role in Pipeline} \\
\midrule
PostgreSQL & 5432 & postgres:16-alpine & Persistent storage for users, sessions, events, sources, and memories \\
Redis & 6379 & redis:7-alpine & Transient state, caching, and queue-assist operations \\
AI Engine & 8000 & Python + FastAPI & Agent execution, provider access, search orchestration, report logic \\
Backend API & 5000 (internal) & Node.js + Express & Auth, session lifecycle, event APIs, workspace management \\
Worker & internal & Node.js  & Asynchronous processing of queued research sessions \\
Frontend & 3000 & Next.js & User interactions, progress visualization, output consumption \\
\bottomrule
\end{tabular}
\end{table}
\subsection{Technology Foundations}

The MARP platform is built on a multi-layered technology stack, where each layer is designed to handle a specific set of responsibilities. This layered architecture ensures modularity, scalability, and efficient system coordination. 

The \textbf{AI layer} manages intelligent reasoning, multi-agent coordination, and model execution. The \textbf{backend layer} handles session management, authentication, and system-level operations. The \textbf{frontend layer} provides an interactive interface for visualization and user interaction, while the \textbf{data layer} ensures reliable storage and retrieval of structured and unstructured information.

\subsubsection*{AI Engine Stack}

The AI engine forms the core of the system, enabling advanced reasoning and orchestration of multiple agents. It is primarily built using \textbf{FastAPI}, \textbf{LangChain}, and \textbf{LangGraph}, which together facilitate API handling, agent workflows, and stateful execution.

To support diverse model integrations, multiple LLM provider SDKs are utilized, including \texttt{langchain-openai}, \texttt{langchain-groq}, \texttt{langchain-google-genai}, and \texttt{langchain-ollama}. 

Additionally, the system incorporates machine learning and research utilities such as \texttt{transformers}, \texttt{torch}, and \texttt{chromadb} for embedding storage and retrieval. External knowledge sources are accessed through tools like \texttt{arxiv}, \texttt{wikipedia}, and \texttt{duckduckgo-search}, enabling comprehensive literature acquisition and enrichment.

\subsubsection*{Backend Stack}

The backend layer is implemented using \textbf{Node.js} with the \textbf{Express} framework, providing a robust and scalable server environment. It integrates \texttt{pg} for PostgreSQL database interactions and \texttt{redis} for caching and real-time performance optimization.

Security and session management are handled using \texttt{jsonwebtoken}, while \texttt{multer} enables efficient file uploads within workspace-scoped environments. Logging and monitoring are managed using \texttt{winston}, ensuring traceability and system reliability.

\subsubsection*{Frontend Stack}

The frontend is developed using \textbf{Next.js 14}, combined with \textbf{React 18} and \textbf{TypeScript 5}, ensuring a modern and type-safe development environment. 

For styling and UI consistency, \textbf{Tailwind CSS} is used, while markdown rendering libraries enable dynamic content visualization. Advanced data visualization is achieved using \textbf{ECharts}, and interactive 3D elements are supported through \textbf{Three.js}.

This frontend architecture enables the creation of rich, dashboard-oriented interfaces that provide real-time insights into the research pipeline. Figure~\ref{fig:frontend_landing} illustrates the overall dashboard design and user interaction flow.

\begin{figure}[H]
  \centering
  \includegraphics[width=0.92\textwidth]{assets/frontend_landingpage.png}
  \caption{Landing page with research initiation controls and project overview.}
  \label{fig:frontend_landing}
\end{figure}

\subsection{Agent Inventory and Functional Categories}

The codebase documents a diverse agent inventory. Table~\ref{tab:agent_categories}
groups them by functional role.

\begin{table}[H]
\centering
\caption{Core agent types in the research pipeline.}
\label{tab:agent_categories}
\begin{tabular}{p{2.0cm}|p{7.0cm}|p{2.6cm}}
\toprule
\textbf{Type} & \textbf{Agents} & \textbf{Function} \\
\midrule
Discovery & Domain Intelligence, Historical Review & Build knowledge map \\
Retrieval & Search, Source Validation & Find and rank sources \\
Synthesis & Gap Analysis, Question Engineering & Extract insights \\
Quality & Technical Verification, Hallucination Detection & Verify accuracy \\
Reporting & Scientific Writing, LaTeX Generation & Generate outputs \\
Interaction & Chatbot, Critique & User engagement \\

\bottomrule
\end{tabular}
\end{table}

\subsection{Database Materials}

The schema defines entities for users, sessions, sources, events, uploads, chats, and memory.
Workspace IDs isolate multi-user contexts. Legacy and new session tables coexist through migrations.

\begin{table}[H]
\centering
\caption{Core database entities and their roles in system persistence.}
\label{tab:db_entities}
\begin{tabular}{p{4.2cm}|p{7.5cm}|p{1.8cm}}
\toprule
\textbf{Table} & \textbf{Purpose} & \textbf{Scope} \\
\midrule
\texttt{users} & Account authentication and user profile metadata & User \\
\texttt{api\_keys} & API access tokens and rate limit counters & User \\
\texttt{research\_logs} & Legacy job records and JSON outputs & User \\
\texttt{research\_sessions} & Active research execution state and progress & Workspace \\
\texttt{workspaces} & Project container for session isolation & Workspace \\
\texttt{workspace\_uploads} & User-uploaded files and embedding status & Workspace \\
\texttt{data\_sources} & Retrieved source metadata and URLs & Session \\
\texttt{execution\_events} & Pipeline stage transitions and logs & Session \\
\texttt{chat\_history} & User queries and agent responses & Session \\
\texttt{user\_memories} & Persistent context for agent reasoning & User \\
\texttt{embeddings\_meta} & Vector chunk references for retrieval & Workspace \\
\bottomrule
\end{tabular}
\end{table}

\subsection{Method for This Report Generation}

This report draws from repository files, schemas, manifests, notes, and assets. Every claim traces to discoverable artifacts.

% ---------------------------------------------------------------
\section{Data Preprocessing}
% ---------------------------------------------------------------

Preprocessing encompasses query refinement, source structuring, and payload assembly. External providers return inconsistent results: varying schemas, quality, relevance. Agents need cleaned inputs.

\subsection{Query Refinement and Result Normalization}

Exact-phrase matching works poorly with verbose or awkwardly-phrased research questions.
The platform generates simplified variants automatically: it removes special characters,
succeeding queries with common suffixes like \textit{overview}. If a strict query returns
no results, the system tries variants. This pragmatism prevents empty result sets.

Each search provider returns results in its own schema. Rather than handle each provider
specially throughout the system, the platform normalizes all results into a common structure:
title, URL, description, source ID, favicon, optional date, optional thumbnail. This
uniformity simplifies downstream processing and reduces provider-specific branching logic.
Figure~\ref{fig:research_pipeline_fastapi} shows this pipeline.

\subsection{Resilience and Event Tracing}

Search calls execute in parallel across providers using thread pools. Each call has a timeout.
If one provider fails or hangs, it doesn't block others. Successful results continue to
processing; failed providers emit errors for logging. Partial failures are acceptable. This
design prevents any single provider from collapsing the entire retrieval.

Duplicate results (those sharing the same URL) get merged. Result counts per query are capped
at sensible limits to manage response times. When queries return zero results, fallback links
ensure users never see a blank state.

Every significant stage transition emits an event: which stage completed, which follows next,
what message to show, and associated metadata. These events persist in the database and stream
to connected frontend clients in real time. This traceability enables both live monitoring
and postmortem analysis of failed runs.

% ---------------------------------------------------------------
\section{Quality Assessment and Source Ranking}
% ---------------------------------------------------------------

In machine learning contexts, feature selection often means choosing which numeric attributes
to include in a model. Here, the problem differs: we must rank and filter sources before agents
synthesize them into research. This process answers a simple question: which retrieved sources
contribute most to research quality?

\subsection{Relevance Assessment}

The platform uses multiple signals to assess source relevance:

\begin{enumerate}[label=\textbf{\alph*.}]
  \item Semantic overlap between the research question and the source title.
  \item Semantic overlap between the research question and the source description.
  \item Confidence in source categorization (is this an academic paper, forum post, blog, etc.?).
  \item Presence of structured metadata like publication dates.
  \item URL-based deduplication confidence (avoiding duplicate results).
\end{enumerate}

These signals guide source prioritization. Sources with high semantic alignment, clear
authoritative provenance, and recent dates rank higher than uncertain, undated,  or
lower-authority equivalents.

\subsection{Multi-stage Quality Filtering}

After retrieval, sources pass through dedicated quality agents. These agents verify technical
accuracy, assess reproducibility, detect hallucinations (claims the source doesn't actually
support), and provide critical review. This quality filtering acts as a gate: weak,
contradictory, or untrustworthy sources get flagged or excluded before synthesis. This reduces
the likelihood of weak evidence contaminating final research outputs.

\subsection{Handling Missing and Zero Values}

Missing metadata (nulls) occur frequently: some sources provide publication dates, others don't.
Some include thumbnail images, others don't. Rather than invent placeholder values, the
platform preserves nulls. This design choice maintains provenance integrity---a missing date
should be visibly missing, not replaced with a guess.

Similarly, when a query returns zero results, that's recorded faithfully. Fallback queries
ensure users don't see empty research states, but the true count remains zero. This accuracy
matters for analysts attempting to understand what worked and what didn't.

% ---------------------------------------------------------------
\section{Model Development and Deployment}
% ---------------------------------------------------------------

This section describes how agent execution works in practice, how the system selects among
model providers, and how the backend coordinates asynchronous work.

\subsection{Agent Execution and Provider Selection}

The AI engine runs agent pipelines asynchronously, tracking each job with a session ID.
Execution is resumable from checkpoints, enabling recovery from failures without restarting
from the beginning. Provider selection uses a runtime factory. The operator chooses a mode
through environment variables: OFFLINE (use local Ollama only), ONLINE (use cloud providers),
HYBRID (try cloud, fall back to local), or HUGGINGFACE (use transformer models locally).
This flexibility accommodates different deployment scenarios: development machines run locally,
production systems use cloud APIs, edge deployments run isolated.
Figure~\ref{fig:fastapi_pipeline} illustrates the request-to-result flow.

\subsection{Backend Coordination and Asynchronous Processing}

The backend API queues research requests. A separate worker process pulls jobs from the queue
and delegates to the AI engine. This separation ensures the API remains responsive: users see
immediate confirmation, while heavy processing happens in the background. Workers monitor the
AI engine's health, retrieve results, and update session state in the database. Every significant
event (agent started, sources found, synthesis complete) gets recorded with timestamps,
severity levels, and detailed context. This creates a queryable audit trail.

Outputs from completed research appear in multiple formats: structured JSON for programmatic
use, Markdown for document authoring, LaTeX for academic typesetting.

\begin{figure}[H]
  \centering
  \includegraphics[width=0.92\textwidth]{assets/research_pipeline_fastapi.png}
  \caption{Query processing flow from input normalization through provider routing to result aggregation.}
  \label{fig:fastapi_pipeline}
\end{figure}

% ---------------------------------------------------------------
\section{Results and System Capabilities}
% ---------------------------------------------------------------

The repository doesn't include formal benchmark suites. But the code, configuration, and
example runs do reveal operational capabilities. This section presents those observations.

\subsection{Operational Readiness Indicators}

Production systems require more than correct outputs. They require reliability infrastructure.
Table~\ref{tab:ops_indicators} lists implementation patterns that signal production maturity.

\begin{table}[H]
\centering
\caption{Production patterns implemented in the system.}
\label{tab:ops_indicators}
\begin{tabular}{p{2.8cm}|p{6.8cm}|p{2.2cm}}
\toprule
\textbf{Pattern} & \textbf{Implementation} & \textbf{Benefit} \\
\midrule
Health checks & Docker Compose probes & Startup ordering \\
Retry logic & \texttt{retry\_count} in database & Fault tolerance \\
Event logging & Indexed \texttt{execution\_events} table & Audit trail \\
Isolation & Workspace IDs in schema & Multi-tenant security \\
Fallbacks & Query variants on failures & Reduces errors \\
Provider modes & ONLINE/OFFLINE/HYBRID & Deployment flexibility \\
Async workers & Separate Node.js process & Non-blocking API \\
\bottomrule
\end{tabular}
\end{table}

\subsection{Runtime Visibility and Control}

To verify basic functionality, we tested the AI engine with a query (\textit{machine learning}).
It returned relevant arXiv papers, confirming the retrieval pipeline works. We note that narrow
or speculative queries naturally return fewer results; this reflects the nature of literature
systems, not a system defect.

The frontend provides real-time visibility into running research tasks. As a job executes,
users see live updates: which stage is active, what sources were found, how synthesis is
proceeding. This transparency is valuable both during development (debugging is easier when
you see intermediate states) and in production (users know their request isn't hanging).

\begin{figure}[H]
  \centering
  \includegraphics[width=0.92\textwidth]{assets/research_live_feed.png}
  \caption{Real-time execution monitoring with stage transitions and event timestamps.}
  \label{fig:research_live_feed}
\end{figure}

In parallel, the source discovery panel (Figure~\ref{fig:research_live_source}) shows results
as retrieval agents find them, making search progress visible.

\begin{figure}[H]
  \centering
  \includegraphics[width=0.92\textwidth]{assets/research_live_source.png}
  \caption{Sources accumulating from parallel retrieval across Web, Wikipedia, and arXiv providers.}
  \label{fig:research_live_source}
\end{figure}

Domain analysis views (Figure~\ref{fig:domain_analysis}) expose intermediate synthesis results,
allowing users to inspect how the system is organizing and interpreting sources.

\begin{figure}[H]
  \centering
  \includegraphics[width=0.92\textwidth]{assets/research_live_domain_analysis.png}
  \caption{Domain extraction and knowledge graph assembly from retrieved sources.}
  \label{fig:domain_analysis}
\end{figure}

The end state of a completed run is shown in Figure~\ref{fig:research_done}, where
generated artifacts are available for downstream reporting and review.

\begin{figure}[H]
  \centering
  \includegraphics[width=0.92\textwidth]{assets/research_done.png}
  \caption{Final research outputs available in JSON, Markdown, and LaTeX formats.}
  \label{fig:research_done}
\end{figure}

Interactive follow-up is supported through chat, shown in
Figure~\ref{fig:user_chat_area}, while API endpoints remain auditable via the backend
OpenAPI/Swagger interface in Figure~\ref{fig:backend_swagger}.

\begin{figure}[H]
  \centering
  \includegraphics[width=0.92\textwidth]{assets/user_chat_area.png}
  \caption{Interactive follow-up interface for clarifying questions and deeper analysis.}
  \label{fig:user_chat_area}
\end{figure}

\begin{figure}[H]
  \centering
  \includegraphics[width=0.92\textwidth]{assets/backend_swagger.png}
  \caption{Swagger API documentation with endpoint schemas, authentication, and rate limits.}
  \label{fig:backend_swagger}
\end{figure}


The schema traces from input to output:

\begin{enumerate}[label=\textbf{\alph*.}]
  \item Sessions preserve topic, status, depth, and timestamps.
  \item Sources preserve domain, URL, count, and metadata.
  \item Events preserve stage transitions and payloads.
  \item Chat and memory preserve user interaction context.
  \item Embedding metadata links sources to vectors for retrieval.
\end{enumerate}

Figure~\ref{fig:schema_architecture} illustrates this persistence design.

\begin{figure}[H]
  \centering
  \includegraphics[width=0.92\textwidth]{assets/schema_architecture.png}
  \caption{Database schema with user accounts, workspace boundaries, sessions, events, sources, and embeddings.}
  \label{fig:schema_architecture}
\end{figure}

\subsection{Performance Metrics}

System performance directly impacts user experience and operational cost. Three key metrics reveal how the platform behaves under realistic research workloads.

Agent execution time (Figure~\ref{fig:agent_exec_times}) demonstrates that specialized agents complete their tasks efficiently. Different agents show varying durations: discovery agents take longer for deep analysis, while validation agents complete quickly. This variability is expected and healthy—it reflects the genuine complexity of each task.

\begin{figure}[H]
  \centering
  \includegraphics[width=0.92\textwidth]{assets/metrics_agent_execution_times.png}
  \caption{Agent execution times across pipeline stages, showing discovery and synthesis taking longer than quick validation checks.}
  \label{fig:agent_exec_times}
\end{figure}

Source discovery performance (Figure~\ref{fig:source_disc_metrics}) shows how quickly the retrieval pipeline finds and normalizes sources from multiple providers. The platform retrieves sources in parallel, and the metric shows both individual provider speeds and aggregate throughput. This parallelism explains why total time remains manageable even with multiple providers.

\begin{figure}[H]
  \centering
  \includegraphics[width=0.92\textwidth]{assets/metrics_source_discovery.png}
  \caption{Source discovery throughput and latency across Web, Wikipedia, and arXiv retrieval providers.}
  \label{fig:source_disc_metrics}
\end{figure}

Search performance (Figure~\ref{fig:search_perf_metrics}) tracks query response times and fallback behavior. The metric reveals that exact matches complete quickly, while fallback queries (when exact matches return zero results) take longer due to the query simplification loop. This tradeoff is acceptable: users prefer slightly slower responses that return results over fast responses that return nothing.

\begin{figure}[H]
  \centering
  \includegraphics[width=0.92\textwidth]{assets/real_metrics_search_performance.png}
  \caption{Search query performance showing exact match vs. fallback query latency and result counts.}
  \label{fig:search_perf_metrics}
\end{figure}

These metrics indicate the platform operates within acceptable bounds. Agent execution times scale linearly with task complexity. Retrieval completes in seconds, not minutes. Search fallback logic prevents empty results without excessive delays. The system behaves predictably, making capacity planning and performance tuning straightforward.

% ---------------------------------------------------------------
\section{Discussion}
% ---------------------------------------------------------------

AI workflow systems commonly fail: unclear boundaries, hidden outputs, minimal monitoring. This system avoids these traps through modular agents with explicit contracts, visible intermediate stages, comprehensive event logging, and rigorous health monitoring.

\subsection{What the Implementation Gets Right}

Five architectural patterns emerge as strengths:

\begin{enumerate}[label=\textbf{\alph*.}]
  \item \textbf{Modular Agent Design:} Specialized agents reduce prompt complexity and make
        reasoning tractable. Each agent owns a specific responsibility rather than attempting
        everything in a single monolithic prompt.
  \item \textbf{Production-Grade Deployment:} Microservices architecture with health checks,
        ordered startup, and shared storage ensure consistent behavior across development,
        testing, and production environments.
  \item \textbf{Transparent Execution:} Event logging creates detailed audit trails. Users
        and developers can see exactly what happened and debug failures from stored records
        rather than guessing.
  \item \textbf{Runtime Flexibility:} Environment variables control provider selection and
        execution mode without code changes. This accommodates different infrastructure
        scenarios without forking the codebase.
  \item \textbf{Multi-Tenant Data Isolation:} Workspace IDs thread through all layers,
        enabling safe multi-user and multi-project operation within a single deployment.
\end{enumerate}

\subsection{Limitations and Challenges}

No system is perfect. Real limitations merit acknowledgment:

\begin{enumerate}[label=\textbf{\alph*.}]
  \item \textbf{Evaluation Gaps:} Outputs and execution traces exist, but no standardized
        test suite covers all agents systematically. Future work should develop benchmark
        datasets.
  \item \textbf{External Variability:} Model and search APIs vary by geography, rate limits,
        response format. Occasional failures and missing fields are inherent to external
        dependencies, not system bugs.
  \item \textbf{Query Sensitivity:} Very narrow research questions naturally yield sparse
        results. This is a feature of academic literature, not a platform defect.
  \item \textbf{Documentation Maintenance:} Assets and references occasionally drift (image
        paths, schema descriptions). Real systems require living documentation discipline.
\end{enumerate}

\subsection{Recommended Future Work}

This implementation provides a solid foundation. Promising extensions include:

\begin{enumerate}[label=\textbf{\alph*.}]
  \item Develop standardized benchmark datasets covering all agent categories (retrieval,
        synthesis, validation) with ground truth labels.
  \item Add provider confidence calibration: measure API accuracy empirically and weight
        results accordingly.
  \item Distinguish network transients from validation failures in retry logic, enabling
        smarter fault handling.
  \item Build a mobile client achieving feature parity with the web interface for
        full-lifecycle research support.
  \item Map generated text back to source documents, creating explicit provenance chains
        from claim to evidence.
\end{enumerate}

% ---------------------------------------------------------------
\section{Conclusion}
% ---------------------------------------------------------------

This report examined a Multi-Agent LLM Research Automation Platform from first principles: its problem domain, architectural approach, and design priorities.

The system orchestrates research through modular agents, coordinates multiple model providers via factory patterns, maintains visibility through event logging, and isolates deployments via workspace scoping.

Research automation requires traceability, recovery mechanisms, real-time visibility, and explicit interfaces. This codebase demonstrates all four, providing a concrete reference for production-grade AI systems.

% ---------------------------------------------------------------
% References
% ---------------------------------------------------------------
\begin{thebibliography}{99}

\bibitem{wooldridge1995}
Wooldridge, M., \& Jennings, N. R.: Intelligent agents: Theory and practice.
\textit{The Knowledge Engineering Review}, 1995, 10(2), 115--152.

\bibitem{jennings1998}
Jennings, N. R., Sycara, K., \& Wooldridge, M.: A roadmap of agent research and
development. \textit{Autonomous Agents and Multi-Agent Systems}, 1998, 1(1), 7--38.

\bibitem{russell2020}
Russell, S., \& Norvig, P.: \textit{Artificial Intelligence: A Modern Approach}
(4th ed.). Pearson, 2020.

\bibitem{vaswani2017}
Vaswani, A., Shazeer, N., Parmar, N., et al.: Attention is all you need.
In \textit{Advances in Neural Information Processing Systems (NeurIPS)}, 2017.

\bibitem{devlin2019}
Devlin, J., Chang, M.-W., Lee, K., \& Toutanova, K.: BERT: Pre-training of deep
bidirectional transformers for language understanding. In \textit{NAACL-HLT}, 2019.

\bibitem{brown2020}
Brown, T. B., Mann, B., Ryder, N., et al.: Language models are few-shot learners.
In \textit{Advances in Neural Information Processing Systems (NeurIPS)}, 2020.

\bibitem{ouyang2022}
Ouyang, L., Wu, J., Jiang, X., et al.: Training language models to follow
instructions with human feedback. In \textit{Advances in Neural Information
Processing Systems (NeurIPS)}, 2022.

\bibitem{openai2023gpt4}
OpenAI: GPT-4 technical report. \textit{arXiv preprint arXiv:2303.08774}, 2023.

\bibitem{touvron2023llama2}
Touvron, H., Martin, L., Stone, K., et al.: Llama 2: Open foundation and
fine-tuned chat models. \textit{arXiv preprint arXiv:2307.09288}, 2023.

\bibitem{wei2022cot}
Wei, J., Wang, X., Schuurmans, D., et al.: Chain-of-thought prompting elicits
reasoning in large language models. In \textit{NeurIPS}, 2022.

\bibitem{yao2023react}
Yao, S., Zhao, J., Yu, D., et al.: ReAct: Synergizing reasoning and acting in
language models. In \textit{International Conference on Learning Representations
(ICLR)}, 2023.

\bibitem{shinn2023reflexion}
Shinn, N., Cassano, F., Labash, B., et al.: Reflexion: Language agents with verbal
reinforcement learning. In \textit{NeurIPS}, 2023.

\bibitem{wang2023selfconsistency}
Wang, X., Wei, J., Schuurmans, D., et al.: Self-consistency improves chain of thought
reasoning in language models. In \textit{ICLR}, 2023.

\bibitem{lewis2020rag}
Lewis, P., Perez, E., Piktus, A., et al.: Retrieval-augmented generation for
knowledge-intensive NLP tasks. In \textit{NeurIPS}, 2020.

\bibitem{guu2020realm}
Guu, K., Lee, K., Tung, Z., et al.: REALM: Retrieval-augmented language model
pre-training. In \textit{ICML}, 2020.

\bibitem{borgeaud2022retro}
Borgeaud, S., Mensch, A., Hoffmann, J., et al.: Improving language models by
retrieving from trillions of tokens. In \textit{ICML}, 2022.

\bibitem{mialon2023augmented}
Mialon, G., Dessi, R., Lomeli, M., et al.: Augmented language models: A survey.
\textit{Transactions on Machine Learning Research (TMLR)}, 2023.

\bibitem{ji2023surveyhallucination}
Ji, Z., Lee, N., Frieske, R., et al.: Survey of hallucination in natural language
generation. \textit{ACM Computing Surveys}, 2023, 55(12), 1--38.

\bibitem{bommasani2021foundation}
Bommasani, R., Hudson, D. A., Adeli, E., et al.: On the opportunities and risks of
foundation models. \textit{arXiv preprint arXiv:2108.07258}, 2021.

\bibitem{bubeck2023sparks}
Bubeck, S., Chandrasekaran, V., Eldan, R., et al.: Sparks of artificial general
intelligence: Early experiments with GPT-4. \textit{arXiv preprint arXiv:2303.12712},
2023.

\bibitem{chen2021evaluating}
Chen, M., Tworek, J., Jun, H., et al.: Evaluating large language models trained on
code. \textit{arXiv preprint arXiv:2107.03374}, 2021.

\bibitem{newman2015microservices}
Newman, S.: \textit{Building Microservices}. O'Reilly Media, 2015.

\bibitem{fielding2000rest}
Fielding, R. T.: \textit{Architectural Styles and the Design of Network-based
Software Architectures}. PhD thesis, University of California, Irvine, 2000.

\bibitem{bernstein2014containers}
Bernstein, D.: Containers and cloud: From LXC to Docker to Kubernetes.
\textit{IEEE Cloud Computing}, 2014, 1(3), 81--84.

\bibitem{page2021prisma}
Page, M. J., McKenzie, J. E., Bossuyt, P. M., et al.: The PRISMA 2020 statement:
an updated guideline for reporting systematic reviews.
\textit{BMJ}, 2021, 372:n71.

\bibitem{kitchenham2007slr}
Kitchenham, B., \& Charters, S.: \textit{Guidelines for Performing Systematic
Literature Reviews in Software Engineering}. EBSE Technical Report EBSE-2007-01,
Keele University and Durham University, 2007.

\bibitem{lamport1994latex}
Lamport, L.: \textit{LaTeX: A Document Preparation System} (2nd ed.).
Addison-Wesley, 1994.

\bibitem{knuth1984texbook}
Knuth, D. E.: \textit{The TeXbook}. Addison-Wesley, 1984.

\end{thebibliography}

\end{document}
